% !PRINTED 391
\weberpage{391}
%<<CONT
minante $D$ dieser Ordnung und zugleich Normenreste einer Primzahl $p$ \gesperrt{in bezug auf diese Ordnung sind, d.\,h. der Zahlen $a$, die der Bedingung}

\weq{1}{N(\omega) \equiv a \pmod{p}}

für eine Zahl $\omega$ der Ordnung $O$ genügen. Es gilt zunächst der Satz:

1. \gesperrt{Wenn $p$ nicht in $D$ aufgeht, so ist jede Zahl $a$ in $Z$ Normenrest.}

Es braucht dies nur bewiesen zu werden für den Fall, daß $a$ ein quadratischer Nichtrest ist, da ja jeder quadratischer Rest zugleich Normenrest ist. % [sic] Gedruckt ist quadratischer.

Es sei also $\mathfrak{p}$ ein in $p$ aufgehendes Primideal des Körpers $\Omega$.

Ist dann $D$ quadratischer Rest von $p$, so ist $\mathfrak{p}$ vom \gesperrt{ersten Grade} (\S\,92), und jede Zahl ist nach dem Modul $\mathfrak{p}$ mit einer rationalen Zahl kongruent.

Da $p$ nicht in $D$ aufgeht, so ist das zu $\mathfrak{p}$ konjugierte Ideal $\mathfrak{p}'$ von $\mathfrak{p}$ verschieden, und wir können eine Zahl $\omega$ in $O$ finden, die den Kongruenzen:

\[
\begin{aligned}
\omega &\equiv a \pmod{\mathfrak{p}},\\
&\equiv 1 \pmod{\mathfrak{p}'}\ ^{1)}
\end{aligned}
\]

genügt (Bd.\,II, \S\,166). Ist dann $\omega'$ zu $\omega$ konjugiert, so ist

\[ \omega \equiv a, \quad \omega' \equiv 1 \pmod{\mathfrak{p}}, \]

und folglich

\[ N(\omega) \equiv a \pmod{\mathfrak{p}}. \]

Ist aber $D$ quadratischer Nichtrest von $p$, so ist $\mathfrak{p} = p$, und $D^{\frac{p-1}{2}} \equiv -1 \pmod{p}$, folglich $\sqrt{D}^{\,p} \equiv -\sqrt{D}$ % Source scan places p outside the radical, as transcribed.
und für jede Zahl in $O$

\weq{2}{\begin{aligned} \omega^p &\equiv \omega' \pmod{p},\\ N(\omega) &\equiv \omega^{p+1} \pmod{p}. \end{aligned}}

Ist nun $\gamma$ eine primitive Wurzel von $p$ im Körper $\Omega$, so ist $c \equiv \gamma^{p+1} \pmod{p}$ eine primitive Wurzel von $p$ im Körper der rationalen Zahlen (Bd.\,II, \S\,167).

Setzen wir also

\weq{3}{a \equiv \gamma^{(p+1)x}, \quad \omega \equiv \gamma^x \pmod{p},}

so ist nach (2) $N(\omega) \equiv a \pmod{p}$ und unser Satz somit bewiesen.

\par\noindent\rule{0.25\linewidth}{0.4pt}\par
\begingroup\footnotesize
$^{1)}$ Setzt man $\omega = a + \xi \pi Q$, wo $\pi$ eine durch $\mathfrak{p}$, aber nicht durch $\mathfrak{p}'$ teilbare Zahl ist, so kann man die Zahl $\xi$ aus der Kongruenz $a + \xi \pi Q \equiv 1 \pmod{\mathfrak{p}'}$ bestimmen, und $\omega$ gehört wegen des Faktors $Q$ zur Ordnung $O$.
\endgroup
